\chapter{RESULTS AND EVALUATION}

\label{ch:results}

This chapter demonstrates the artifact and then reports what was measured. It is deliberately
organised around the question each piece of evidence answers, not around the size of the test suite.
Where a number comes from a single machine, a single fixture set or a single observational sample,
it is reported with that limitation attached.

\section{System Demonstration}

The application is a single-page web client served alongside an Express API and a PostgreSQL
database. All three run from one repository, and a development launcher builds the shared contract
package before starting the watchers so that the client and the server never disagree about a
schema.

\subsection{The Editable Curriculum Map}

An account with no program context sees the curriculum as semester columns. Each card carries its
code, name, credits and requirement category, and its state is expressed in text as well as
colour: completed, planned, locked, recommended, or a highlighted prerequisite of a course the
reader is inspecting.

\begin{figure}[h]
\centering
\makebox[\textwidth][c]{\includegraphics[width=1.15\textwidth]{semester-grid.png}}
\caption{Semester-column curriculum map with completion and planning states}
\end{figure}

Selecting a locked course expands the requirement it is waiting on, and hovering it highlights the
chain of courses that must be completed first. The disclosure is a planning aid: it shows what the
declared prerequisites are, not what the institution will offer next term.

\begin{figure}[h]
\centering
\makebox[\textwidth][c]{\includegraphics[width=1.15\textwidth]{prerequisite-highlight.png}}
\caption{A locked course expanded to show the prerequisite it is waiting on}
\end{figure}

Elective groups open in a sidebar that lists the group's members and their states. A completed
course is claimed by the group where it was clicked and is dimmed in the other groups that also
list it, so it is counted once.

\begin{figure}[h]
\centering
\makebox[\textwidth][c]{\includegraphics[width=1.15\textwidth]{elective-sidebar.png}}
\caption{Elective group sidebar showing group members and claimed courses}
\end{figure}

\subsection{Context-Aware Screens}

When a reference program is attached to an account, the curriculum route becomes a read-only
reference view. It lists the program's placements with their prerequisite edges, shows each
course's difficulty estimate together with the number of votes behind it, names the requirements
the import could not resolve, and separates current member progress from records the account earned
outside the program. Nothing on that screen is editable, and the screen states why: editing inside
an assigned program requires a claim-transfer policy and a context-keyed client cache that are not
part of this artifact.

Terminal output of the assigned screen is shown below: the reference curriculum, the account's
current member progress, and the saved course list, all read-only.

\begin{figure}[h]
\centering
\makebox[\textwidth][c]{\includegraphics[width=1.15\textwidth]{reference-view.png}}
\caption{The curriculum route for an account with an assigned reference curriculum}
\end{figure}

The same account's totals are honest about their own scope: ten completed courses and thirty-four
earned credits, with the percentage and remaining requirements explicitly withheld because the
program's own rules are unverified.

The planner follows the same branch. A guest receives the legacy workload analysis; an assigned
account receives a context preview that resolves prerequisites, credits and the rating prior inside
that program, rejects a non-member selection, and reports courses it could not schedule with the
reason.

\begin{figure}[h]
\centering
\makebox[\textwidth][c]{\includegraphics[width=1.15\textwidth]{planner-preview.png}}
\caption{The contextual semester preview for the same account}
\end{figure}

That preview shows three things a guest never sees: the grade-based branch resolved from the
account's own average, the difficulty prior drawn from ratings inside the program rather than
globally, and the requirement the import could not resolve, held apart from the course list instead
of being guessed at. The grades screen reports the credit-weighted average, the retake history and the
highest-score distribution, plus a per-semester table when attempts carry term metadata. The ratings
screen lists only completed courses and holds a submitted vote's retry identity until the server
confirms it.

\subsection{Administrative Simulation}

The administrative surface presents a scenario: a reference program, a semester and a year. Within
that scenario an administrator saves resource settings, reads a planned-selection report, compares
declared seats with current demand, previews a capacity envelope, runs either allocation model,
captures the result and browses the history of earlier captures. Every screen in that surface
carries the same qualification: the numbers are simulation assumptions, not university capacity,
and a capture is a record of a simulation rather than a registration decision.

\section{Verification Evidence}

\subsection{Scale and Shape of the Suite}

At the time of writing the automated suite comprises three thousand and seventy-six server tests across
one hundred and nineteen suites, two thousand five hundred and sixteen client tests across one
hundred and ten files, and ten startup lifecycle cases that exercise the real development launcher, including
shutdown of all three watchers from one signal.

Counting tests is not itself evidence, so the suite is better described by what it protects.

\begin{table}[h]
\centering
\small
\begin{tabular}{p{0.32\textwidth}p{0.58\textwidth}}
\toprule
\textbf{Property} & \textbf{Representative evidence} \\
\midrule
Prerequisite enforcement & completion refused on an unmet mandatory parent; recommendations never
offer a blocked course \\
Transitive cascade & uncompletion removes direct and transitive dependents in one transaction,
including a real concurrent writer \\
Scope isolation & a mismatched expected scope conflicts and leaves stored rows byte-identical;
cascades refuse across owners \\
Immutable attempts & identical retry returns the original row; changed payload under a used key is
refused; boundaries 0 and 100 enforced at the database \\
Rollback & a real failure after each write stage leaves no partial run, participant or receipt \\
Retry recovery & a lost reply is recovered by an exact retry; a changed payload conflicts \\
Corruption & malformed or mismatched stored history fails closed rather than displaying a partial
page \\
Honest absence & unknown credits, seats and eligibility produce their own states, never zeros \\
Access control & unauthenticated, non-owner, non-administrator and demoted-mid-transaction calls
are refused \\
\bottomrule
\end{tabular}
\caption{What the suite is organised to protect}
\end{table}

\subsection{Fail-First Practice}

Each increment's new tests were run against the code they replace before the fix was applied, and
the failure was recorded. This matters because a test written after a fix tends to encode the
implementation rather than the requirement. A representative case is the session-expiry fix, whose
seventeen integration and five unit tests were each observed failing without the change.

The same discipline applies to defects found during review. Where a review identified a gap -- an
impossible prerequisite that ranking still rewarded, a numeric precision edge in a stored JSON
position, a missing independent guard on a profiling helper -- the test was written to fail against
the pre-fix source and then kept.

\subsection{Reproducibility}

The runtime is pinned to an exact Node and npm pair, the dependency lockfile is committed, and
installs use the lockfile rather than resolving ranges. One continuous-integration job runs the
whole gate on every change: restore, build the shared package, build both applications, typecheck,
lint with zero warnings, check formatting, apply migrations to a disposable PostgreSQL service,
seed it, then run the startup, server and client suites. The job never touches an application
database and holds no deployment secret, so a contributor cannot mistake it for a deployment
pipeline.

\subsection{An Open Finding}

Full server runs against the same seeded disposable database have failed different suites, and
every failing suite has passed in isolation. One run produced a 401 where a business-level refusal
was expected. A later run failed two suites untouched by the change under test, both returning a
conflict where success was expected. Chasing the second case to its source identified the mechanism:
that status is not a business refusal but the final throw after a bounded serializable retry loop
exhausts every attempt, a path both services share. The assertions therefore test that four
concurrent writers all succeed against a loop that may legitimately give up under contention, and a
loaded runner supplies it. The retry bound was left unchanged, because one contention-sensitive run
is not evidence for widening it; the open question is why every attempt collides, not how many to
permit. Recording a mechanism rather than a symptom is the difference between a failure that costs
an hour to re-derive and one that is recognised on sight.

\section{Measured Behaviour of the Simulation Models}

The allocation models are the only parts of the artifact whose cost is not dominated by network or
database latency, so they were measured directly with a purpose-built profiler over synthetic
fixtures at the supported bounds: up to five hundred students, one hundred choices per student, ten
thousand catalogue entries and one hundred rounds.

\begin{table}[h]
\centering
\small
\begin{tabular}{lrr}
\toprule
\textbf{Scenario} & \textbf{Allocation median (ms)} & \textbf{Additional replay median (ms)} \\
\midrule
Maximum students and choices & 3381--4843 & 1681--2502 \\
\bottomrule
\end{tabular}
\caption{Local observations of the semester allocation model}
\end{table}

The observations were taken on an Apple M2 running macOS on arm64, with one warm-up and three
samples per scenario. Full result digests agree across repeats and across reversed input order, and
input digests confirm the fixtures were not mutated, so the measurement is repeatable even though it
is local.

A second profiler measured the persisted verification and projection paths that a stored capture
must replay. Its maximum single-run observations were between 1.7 and 2.5 seconds for one capture,
and between 10.4 and 17.7 seconds for a history page, using one observational sample per scenario.

Three limitations attach to all of these numbers. They exclude database time, capture time, worker
time and HTTP latency; they come from one machine; and the history figure in particular rests on a
single sample rather than a distribution. They therefore bound nothing about deployed latency, and
they are reported as observations about the model rather than as performance guarantees.

\section{Accessibility Evaluation}

Because the interface is dense, correctness of rendering was measured rather than assumed. The audit
ran in headless Chromium with computed styles and an automated rule engine, at desktop and mobile
widths, with scripted keyboard traversal.

The audit produced nine defects, all of which are fixed: a map that faded when a locked course was
hovered; an elective panel more than forty tab stops from its trigger and indifferent to Escape; a
save that lost keyboard focus because the focused control became disabled; zoom controls that looked
available at their limits; a crash screen with no exit; a sign-in error that left focus at the top of
the document; a rejected score that did not identify the offending field; a curriculum load that
could not be retried in place; and the session-expiry race.

One finding was deliberately not fixed. At a 390-pixel viewport the rule engine reports a
label-in-name mismatch on curriculum cards, because the card extends past the edge of a horizontally
scrollable region, so the engine drops the credits from the visible text and compares the remainder
with the accessible name. The accessible name does contain every visible word and the page is clean
at 1280 pixels, so the name was left unchanged rather than distorted to satisfy the tool. The
disposition is recorded with its reasoning.

\section{Threats to Validity}

\textbf{Construct validity.} The artifact measures internal consistency, not academic truth. A
credit total that the program does not declare is stored as unknown, and no screen presents a
derived graduation total while the program's rules remain unverified. A reader could reasonably
want a number; the artifact's position is that a number would be wrong.

\textbf{Internal validity.} The simulated allocation models are greedy, scarcity-ordered heuristics.
They are not optimal, and no optimality claim is made. Their outputs are reproducible because the
contracts are frozen and the tie-breaking is deterministic, which makes them inspectable, not
correct.

\textbf{External validity.} Four program structures are stored, three of them transcribed from a
signed document: Data Science at 124 credits and the two Information Technology specialisations at
150 each. All three reconcile, and all three are imported as reference structures that assign no
account; the reference program is the author's own reading of the current published page. The earlier
expectation that the other programs' sources would contradict themselves proved wrong in a specific
way: each document disagreed with itself in places, and every such disagreement was resolved by
reading the rows rather than the labels. Where a document did not state something the transcription
needed, the manifest records the rule it applied instead of presenting an inference as a quotation.
Breadth across programs is therefore demonstrated; breadth across schools is not.

\textbf{Ecological validity.} No students were recruited, no planning behaviour was observed, and
no claim is made that the recommendations improve outcomes. The evidence is that the artifact
enforces the invariants it declares, is reproducible, and remains operable by keyboard and screen
reader at the widths tested.

\textbf{Construct of the measurements.} The profiling figures are local observations of pure
functions on synthetic data. They say something about the models and nothing about the deployed
system, and the report states the excluded costs alongside each number.
